\documentclass[12pt]{article}
\usepackage[a4paper,margin=18mm]{geometry}
\usepackage[T1]{fontenc}
\usepackage{amsmath,amssymb,amsthm}
\usepackage{enumitem}
\usepackage{hyperref}
\usepackage[nameinlink]{cleveref}

\setlength{\parindent}{0pt}
\pagestyle{empty}

\theoremstyle{definition}
\newtheorem{definition}{Definition}[subsection]

\begin{document}

\begin{center}
  {\LARGE\bfseries Lecture 1: Real Numbers}\\[0.25em]
  {\large AMA3707 Real Analysis}\\[0.15em]
\end{center}

\section{Real Numbers}
\subsection{Numbers and Fields}

\begin{definition}[Natural Numbers]
\label{def:natural-numbers}
The \emph{natural numbers} are the positive integers:
\[
  \mathbb{N}=\{1,2,3,\ldots\}.
\]
In these notes, \(0\notin\mathbb{N}\).
\end{definition}

\begin{definition}[Integers]
\label{def:integers}
The \emph{integers} are the whole numbers, including the positive integers,
the negative integers, and zero:
\[
  \mathbb{Z}=\{\ldots,-2,-1,0,1,2,\ldots\}.
\]
\end{definition}

\begin{definition}[Rational Numbers]
\label{def:rational-numbers}
A real number is called \emph{rational} if it can be expressed as a quotient
of two integers. The set of rational numbers is
\[
  \mathbb{Q}
  =\left\{\frac{p}{q}:p,q\in\mathbb{Z},\ q\neq0\right\}.
\]
\end{definition}

\end{document}

